% Generated draft: requires mathematical and source review.
\documentclass[11pt]{article}
\usepackage[margin=25mm]{geometry}
\usepackage{fontspec}
\setmainfont{Latin Modern Roman}
\usepackage{amsmath,amssymb,mathtools,unicode-math}
\setmathfont{Latin Modern Math}
\usepackage{enumitem,needspace,xurl,hyperref}
\setlength{\parindent}{0pt}
\setlength{\parskip}{5pt}
\setlength{\emergencystretch}{3em}
\newcommand{\N}{\mathbb{N}}
\newcommand{\Z}{\mathbb{Z}}
\newcommand{\Q}{\mathbb{Q}}
\newcommand{\R}{\mathbb{R}}
\newcommand{\C}{\mathbb{C}}
\newcommand{\F}{\mathbb{F}}
\newcommand{\A}{\mathbb{A}}
\newcommand{\PP}{\mathbb{P}}
\newcommand{\E}{\mathbb{E}}
\newcommand{\eps}{\varepsilon}
\newcommand{\dd}{\,\mathrm{d}}
\begin{document}
\section*{UnsolvedMath ID: 8800063}
\textbf{Source problem code:} catalogue entry
\subsection{Definitions and mathematical statement}
For a full-rank Euclidean lattice \(L\subset\R^n\), let \(\lambda_j(L)\) be the smallest radius containing \(j\) linearly independent lattice vectors. The \(\gamma(n)\)-approximate shortest-vector problem asks for \(v\in L\setminus\{0\}\) with
\[
\|v\|_2\le\gamma(n)\lambda_1(L).
\]
The corresponding shortest-independent-vectors problem asks for linearly independent \(v_1,\ldots,v_n\in L\) with
\[
\max_i\|v_i\|_2\le\gamma(n)\lambda_n(L).
\]
An ideal-lattice instance is a full-rank ideal of a specified rank-\(n\) ring, embedded into Euclidean space using a specified coefficient or canonical embedding. The ring family and embedding are part of the computational problem; arbitrary cyclic lattices and ideals in number-field orders are not automatically the same class.

Micciancio asks for an analogous small-polynomial-factor worst-case hardness result for these two problems on ideal lattices. A polynomial factor would have the form \(\gamma(n)\le Cn^c\) for fixed constants, but the intended small exponent and baseline theorem are not supplied. A constant-factor result is a different, stronger approximation requirement and is not substituted for the stated polynomial-factor target.

The word analogous refers to missing earlier material. The source does not specify whether hardness means NP-hardness under a chosen reduction, a reduction from general lattices, or a worst-case-to-average-case reduction underlying a cryptographic assumption. These are different statements with different directions and parameters. No claim of one type is invented from another.

The supplied research summary mixes NP-hardness claims, approximation factors described as subexponential/superpolynomial, and a remaining constant-factor question, without exact theorem statements. Under a fixed NP-hardness notion, hardness for an easier large-factor approximation would already imply hardness for the harder smaller-factor search task, so these phrases cannot simply be read as an ordered list of NP-hardness bounds. The exact reductions and ring restrictions must be recovered before assigning a resolved scope.
\subsection{Short English statement}
Establish the intended worst-case hardness theorem for polynomial-factor approximate SVP and SIVP on a specified ideal-lattice family. The antecedent theorem, reduction notion, ring embedding, and factor exponent are missing, and the supplied hardness summary mixes distinct notions.
\subsection{Sources}
\begin{itemize}
\item Local UnsolvedMath registry, permanent ID \(8800063\), exact source code \texttt{catalogue entry}, ``Polynomial-factor hardness of ideal-lattice problems.'' Proposer Daniele Micciancio.
\item IPAM \emph{Workshop on Number Theory and Cryptography --- Open Problems} (2006), attachment 6240, recorded PDF page 27; \url{https://www.ipam.ucla.edu/programs/workshops/workshop-i-number-theory-and-cryptography-open-problems/?tab=schedule}. Recorded access 2026-07-29; pdf-and-legacy-ppt visual-audit extraction.
\item Supplied PARTIAL-PROGRESS summary names Micciancio (2002), Peikert--Rosen, Lyubashevsky--Micciancio, and an unspecified 2013 sharpening, without complete references or precise hardness/reduction statements.
\end{itemize}
\end{document}
